\chapter{Monetary Policy}

Central banking enters this lecture in two stages. We begin with the familiar modern description: a central bank is the institution responsible for currency and money. But we very quickly step back, because the point of the lecture is that central banks were not handed down as an obvious part of government. They were invented. A workable design appeared, it solved important problems, and then it was copied. From that starting point the lecture moves from the Bank of England to the Federal Reserve, from reserve requirements to capital requirements, and from older money-multiplier stories to the balance-sheet logic that dominated the recent financial crisis.

\section{Central Banks as Financial Inventions}

A central bank is introduced first in its modern form: a special bank responsible for the currency of a country. But that description is only a surface description. To understand what the institution is, we go back to the older private banking system out of which it emerged.

The lecture recalls the goldsmith bankers. A depositor left gold with a banker and received a paper claim. Those claims then began to circulate. In that sense paper money starts not as an act of the state, but as a private contractual device. The arrangement can work for a long time, but it has a permanent weakness: if too many claimants come back at once, the banker may not be able to redeem the paper in specie.

That weakness motivates the turn to the Bank of England in 1694. The first point is scale. It had a special charter, it was large, and it could act on the rest of the banking system. The second point is more important than the first: a large bank that can destroy small banks by demanding redemption may instead choose a stabilizing role. That is the institutional move that matters.

\subsection{Question \& Answer}

\paragraph{Question.} Why does forcing smaller banks to keep deposits with a dominant bank stabilize the system?

\paragraph{Answer.} The dominant bank can always threaten instability. If it presents another bank's liabilities for payment all at once, the weaker bank may fail. But if the weaker bank keeps a deposit with the dominant bank, the relation changes. The deposit serves as a reserve pool and also as a channel of support. The same institution that could destroy the smaller bank can now lend to it or release funds to it in a crisis. A coercive power becomes a stabilizing power. That is the first central-banking mechanism in the lecture.

This is why the Bank of England matters here not merely as a historical name, but as a model. In the lecture it becomes the template for central banking more generally: reserves at the central institution, support in trouble, and a banking system made more stable because the strongest bank has chosen a ``live and let live'' policy.

\section{The American Path to the Federal Reserve}

The lecture then turns to the American case. Here the point is not that the United States produced something wholly new. The point is almost the opposite: the United States kept circling toward a Bank of England type solution, while often resisting the admission that this was what it was doing.

The Suffolk Bank in Boston is an early example. It was a private bank, but it required New England banks to keep deposits with it, and by doing so it stabilized the region. The lecture uses this example to reinforce the idea that central banking began as an institutional innovation before it became a fully governmental one.

The disorder of nineteenth-century American currency makes the need vivid. Different bank notes traded at different discounts. Merchants had to inspect the issuing bank and apply a discount schedule. The National Banking Acts of 1863 and 1864 addressed that problem by creating national banks whose notes were backed through deposits with the Treasury. That was a major success for currency uniformity: the paper-money problem was largely solved.

But the lecture insists on a distinction that must not be lost. Fixing the currency is not the same as fixing the banking system. Even after national bank notes traded at par, there were still bank runs, banking crises, and credit contractions. In the lecture this separation of problems is decisive, because it explains why the United States eventually moved toward the Federal Reserve.

The Federal Reserve System, created in 1913 and opened in 1914, is then presented as an American adaptation of the same stabilizing logic. It looks different. There are twelve regional reserve banks and a board in Washington. But the institutional mechanics are familiar: member banks hold reserves either as currency in vault or as deposits with the Reserve Banks, and when they get into trouble the Federal Reserve can lend to them. That is why the Fed becomes the lender of last resort.

The discount window belongs here as the concrete operational version of that role. A troubled bank does not simply ask for money. It brings collateral. The central bank discounts the collateral and lends against it. The lecture stresses the old-fashioned institutional language because it captures the nature of the transaction: this is support, but not gratuitous support.

\section{From Bank Rescue to Macroeconomic Stabilization}

The lecture does not let the founding of the Federal Reserve stand as a triumphal ending. There is an immediate reversal. The system looked successful, but the crisis of the early 1930s showed that the Federal Reserve had not put banking panics behind us forever. Banks failed after 1929, the banking holiday arrived in 1933, and Roosevelt shut down the system altogether. The lecture adds federal deposit insurance as the next layer of institutional repair.

From here the mission broadens. The central bank is no longer only the institution that stands behind the banking system in panic. It becomes the institution that is meant to stabilize the whole economy. This is where William McChesney Martin's phrase about taking away the punch bowl enters. The central bank restrains booms before they become destructive, and it eases when the economy becomes weak.

The same theme appears in the lecture's steam-engine metaphor. A regulator is not the engine itself. It is the device that prevents the system from running too hot. So too the Federal Reserve becomes, in the lecture's language, a regulator for the entire economy.

The lecture also pauses over institutional credibility. The European Central Bank is mentioned as a newer example of the same general design, and central-bank independence is emphasized as a response to the fear that ordinary politics may prefer inflationary short-run expedients. In this logic, long terms and insulation from dismissal are not decorative constitutional details. They are part of how a currency becomes stable.

The modern operational center of the discussion is the Federal Open Market Committee. It sets a target range for the federal funds rate, the overnight unsecured rate in lending between banks and certain other financial institutions. At the lecture date the numbers are:

\begin{align}
0 \le i_{ff} \le 0.25\%, \qquad i_{ff} \approx 0.13\%.
\end{align}

The policy interpretation is straightforward. When unemployment is high and the economy is weak, the Federal Reserve pushes this rate toward zero. When inflationary pressure is building, it moves in the other direction.

A second modern operating rate enters after 2008, when the Fed begins to pay interest on reserve balances:

\begin{equation}
i_{or} = 0.25\%.
\end{equation}

The lecture treats this as a genuine regime change. The old operating procedure worked indirectly through open-market operations in Treasury bills and through the interlocking of money markets. Interest on reserves adds a new instrument.

\subsection{Question \& Answer}

\paragraph{Question.} Why can the interest on reserves exceed the federal funds rate?

\paragraph{Answer.} If a member bank is eligible to receive interest on reserves, and if
\[
i_{or} > i_{ff},
\]
then that bank has no reason to lend funds in the federal funds market rather than leave them at the Federal Reserve. The lecture's answer is that the fed funds market does not disappear because some lenders in that market are not eligible to earn interest on reserves. In particular, government-sponsored enterprises can continue to lend there. So the two rates need not coincide, and this divergence is itself evidence that the operating regime has changed.

\section{Reserve Requirements Under Regulation D}

After the discussion of modern operating tools, the lecture narrows again. We return to reserve requirements because they were once the central textbook mechanism of monetary control, and because one cannot understand the later turn to capital requirements without first reconstructing this older framework.

In the lecture's language, reserves are the sum of vault cash and deposits at the Federal Reserve. In symbols we may write
\begin{equation}
R = C_{\text{vault}} + D_{\text{Fed}}.
\end{equation}
The lecture describes these as ``high-powered money'' from the standpoint of the banking system.

Regulation D then enters as the rule that ties reserves to liabilities. At the lecture date the key visible comparison is between transaction accounts and time deposits. The two board images belong here because they preserve exactly this classification.

\begin{figure}
\centering
\includegraphics[width=0.82\textwidth]{lecture_18_figure_02.png}
\caption{Reserve requirements under Regulation D. The board sets up the liability categories before the numerical contrast is fully sharpened.}
\end{figure}

\begin{figure}
\centering

\includegraphics[width=0.82\textwidth]{lecture_18_figure_03.png}
\caption{Transaction accounts versus time deposits. The numerical contrast is clearer here, and the lecture is already moving toward the money-multiplier discussion.}
\end{figure}

The cleaned comparison is:

\begin{align}
\text{Transaction accounts:} &\quad 10\%, \\
\text{Time deposits:} &\quad 0.
\end{align}

Or, in compact notation,
\begin{align}
rr_{\text{TA}} &= 10\%, \\
rr_{\text{TD}} &= 0.
\end{align}

Transaction accounts are the liabilities that can be used for spending and withdrawn immediately. Time deposits are different in the lecture's framing: they are not due on demand in the same way, and so they do not carry the same run risk.

\subsection{Question \& Answer}

\paragraph{Question.} Why are transaction accounts subject to reserve requirements while time deposits are not?

\paragraph{Answer.} The logic is older than any particular regulation. A transaction account can be demanded back at once. If many depositors demand it simultaneously, the bank may face a run. A time deposit does not create that same immediate pressure, because repayment can be delayed under the terms of the contract. Reserve requirements are therefore tied to the liabilities that generate the classic run problem. That is why the lecture treats reserve requirements as a bank-run technology.

From here the lecture reconstructs the older money-multiplier argument. If the reserve ratio is \(rr\), and if banks hold only the reserves they are required to hold, then the deposit base can expand until reserves are the required fraction of deposits. In its simplest form the lecture writes this as
\begin{equation}
m \approx \frac{1}{rr}.
\end{equation}
With a \(10\%\) reserve requirement, this gives the familiar simplification
\begin{equation}
m \approx 10.
\end{equation}

\begin{remark}
The lecture is explicit that this is an oversimplification. It is a useful first model, not an exact identity. The point is to recover the old textbook logic before asking whether it still governs the present regime.
\end{remark}

Under that old theory, if banks merely meet reserve requirements and do not hold excess reserves, then reserves pin down the total quantity of deposits. This is why reserve requirements once seemed to be the central lever of monetary control.

But the lecture then overturns the classical emphasis. Banks now hold large excess reserves. At the lecture date the quantity is described as enormous, on the order of
\begin{equation}
\text{Excess reserves} \gtrsim 1.2 \text{ trillion dollars}.
\end{equation}
Once the banking system is willing to sit on reserves far above the legal minimum, reserve requirements cease to bind in the same way. The lecture's phrase is that ``the world has changed.'' That is the bridge to capital requirements.

\section{Capital Requirements and Risk-Weighted Assets}

The next move is conceptually sharp. Reserve requirements are tied to liabilities. Capital requirements are not. Reserve requirements say: hold reserves against particular liabilities. Capital requirements say: hold common equity against risky assets.

The lecture uses Basel III as the modern framework. In cleaned notation we may summarize risk-weighted assets as
\begin{equation}
\text{RWA} = \sum_i w_i A_i,
\end{equation}
where \(A_i\) is an asset category and \(w_i\) is the associated regulatory risk weight.

Cash carries zero weight in the lecture's example:
\begin{equation}
w_{\text{cash}} = 0.
\end{equation}
Corporate loans carry full weight:
\begin{equation}
w_{\text{corp}} = 1.
\end{equation}

The common-equity requirements then take the form
\begin{align}
\frac{E_{\text{common}}}{\text{RWA}} &\ge 4.5\%, \\
\frac{E_{\text{common}}}{\text{RWA}} &\ge 7\% \quad \text{with the conservation buffer}, \\
\frac{E_{\text{common}}}{\text{RWA}} &\ge 9.5\% \quad \text{with the additional countercyclical buffer.}
\end{align}

\begin{remark}
The spoken lecture is somewhat loose in its verbal buffer arithmetic, but the intended structure is clear: a \(4.5\%\) common-equity minimum, a \(2.5\%\) conservation buffer, and a further buffer that can be added in advance of a bubble.
\end{remark}

\subsection{Question \& Answer}

\paragraph{Question.} What is the difference between reserve requirements and capital requirements?

\paragraph{Answer.} Reserve requirements are a liquidity rule tied to certain liabilities. They ask whether the bank holds enough immediately available resources against claims that may run. Capital requirements are a solvency rule tied to the asset side. They ask whether the bank has enough loss-absorbing common equity relative to the riskiness of what it owns. A bank may satisfy one and fail the other. The lecture's worked example is designed precisely to show this.

The lecture then adds the complication of credit ratings. For many years regulatory systems treated ratings by agencies such as Moody's and Standard \& Poor's as if they were authoritative measurements of risk. Basel III still reflects some of that history. But after the financial crisis, the United States moved away from formal regulatory reliance on ratings. This creates a tension: modern capital regulation wants a way to measure risk, but the obvious rating system has lost much of its credibility.

\section{A Worked Example: A Bank Balance Sheet}

The lecture's mathematical payoff is a miniature bank. We begin with the simplest possible balance sheet identity:
\begin{equation}
A = L + E,
\end{equation}
where \(A\) denotes assets, \(L\) liabilities other than equity, and \(E\) common equity.

\subsection{Initial deposit, reserves, and missing equity}

Suppose a new bank opens its doors and a depositor brings in \$100. We may write the balance sheet as
\begin{align}
A &= 100 \text{ cash}, \\
L &= 100 \text{ transaction account}, \\
E &= 0.
\end{align}

If the deposit is a transaction account, required reserves are
\begin{equation}
0.10 \times 100 = 10.
\end{equation}
The bank holds \$100 in cash, so reserve compliance is easy.

Now look at the asset weights. Since cash carries zero risk weight,
\begin{equation}
\text{RWA} = 0.
\end{equation}
This immediately exposes the distinction between the two regulatory systems. The bank is liquid enough from the standpoint of reserves, but it has no common equity at all. So it is not a meaningful capitalized institution.

The next move is to issue shares. Suppose the bank raises \$20 in common equity. Then the balance sheet becomes
\begin{align}
A &= 120 \text{ cash}, \\
L &= 100 \text{ transaction account}, \\
E &= 20 \text{ common equity}.
\end{align}
Because all assets are still cash, \(\text{RWA}=0\). So the bank now satisfies both reserve and capital requirements.

\subsection{Moving into risky assets}

A bank that holds only cash makes no money. So the lecture next converts \$100 of cash into corporate loans while keeping \$20 in cash:
\begin{align}
A &= 100 \text{ corporate loans} + 20 \text{ cash}, \\
L &= 100 \text{ transaction account}, \\
E &= 20.
\end{align}
Now the risk-weighted assets are
\begin{equation}
\text{RWA} = 100.
\end{equation}
At a \(7\%\) common-equity requirement, the minimum equity is
\begin{equation}
0.07 \times 100 = 7.
\end{equation}
The bank has \(20\), so capital compliance still holds.

Reserve compliance also holds. Against the \$100 transaction account the bank needs only \$10 in reserves, and it is holding \$20 in cash. So the bank is both liquid enough and capitalized enough.

\subsection{Losses, equity wipeout, and the regulatory split}

Now the crisis arrives. Suppose \(20\%\) of the corporate-loan book defaults. Then the loans are written down from \$100 to \$80:
\begin{align}
A &= 80 \text{ corporate loans} + 20 \text{ cash} = 100.
\end{align}
The deposit liability is still \$100. Therefore the balance sheet identity forces common equity to zero:
\begin{align}
L &= 100 \text{ transaction account}, \\
E &= 0.
\end{align}

This is the critical moment in the lecture. Reserve compliance is still satisfied. The bank still has \$20 cash, and the reserve requirement against the \$100 transaction account is only \$10. But capital compliance fails completely. The bank has no common equity left.

We therefore see explicitly that reserve requirements and capital requirements do different jobs. A bank can have enough reserves to meet withdrawals and yet be insolvent or nearly insolvent from the standpoint of equity.

\subsection{Question \& Answer}

\paragraph{Question.} Why can capital requirements become procyclical and dangerous in a crisis even if they are supposed to make banks safer?

\paragraph{Answer.} Once losses wipe out equity, the bank must repair its capital position. In principle it can issue new shares or sell risky assets. But in a crisis both routes are bad. New equity is hard to raise because investors distrust weak banks. Asset sales are also hard because every bank is trying to sell at the same time, so prices fall and balance sheets deteriorate further. What was designed as a stabilizing rule can therefore intensify the downturn when all institutions move together.

The lecture dramatizes this by briefly considering an asset sale and then correcting the idea. If common equity has already fallen all the way to zero, merely shifting the composition of assets from loans to cash does not create new equity. One must actually raise fresh capital.

That is also why the lecture returns, at the end, to the lender of last resort. In the recent crisis the central banks of the world prevented system-wide collapse by lending aggressively when private balance-sheet repair became impossible. Basel III's countercyclical buffer is an attempt to move the adjustment earlier in time, by forcing thicker cushions before the crash rather than only during it. The Dodd--Frank restrictions on discretionary rescue then make the design problem still sharper: if emergency judgment is constrained, ex ante capital regulation has to do more of the work.

\section{Summary}

The chapter has moved in the same order as the lecture. Central banks begin as copied financial inventions, not as obvious bureaucratic appendages of the state. The Bank of England supplies the first clear model: concentrated reserves, systemic oversight, and last-resort support. The Federal Reserve is then an American version of the same logic, later enlarged into a macroeconomic regulator through the federal funds rate and related tools.

The mathematical heart of the lecture lies in the transition from reserve requirements to capital requirements. Reserve requirements once supported the familiar money-multiplier story, but that story weakens once banks willingly hold large excess reserves. Capital requirements then move to the center. The worked balance-sheet example shows why. Reserves protect against runs on demand liabilities; common equity absorbs losses on risky assets. A bank may satisfy the first and fail the second. And once many banks fail the second together, the attempt to restore capital can itself become destabilizing. That is why the role of central banks and regulators remains, in the lecture's final mood, unfinished and still evolving.